% Part: computability
% Chapter: computability-theory
% Section: fixed-point-thm

\documentclass[../../../include/open-logic-section]{subfiles}

\begin{document}

\olfileid[fr]{cmp}{thy}{fix}
\olsection{Le théorème du point fixe}

Revenons au problème de l'arrêt. Employons
provisoirement la notation
$\gn{\cfind{x}(y)}$ pour $\tuple{x,y}$,
que l'on peut voir comme un « nom » de la valeur
$\cfind{x}(y)$. Elle permet de reformuler
l'une de nos preuves de l'indécidabilité du
problème de l'arrêt.

Existe-t-il une fonction calculable~$h$
ayant la propriété suivante ? Pour tous
$x$ et~$y$,
\[
h(\gn{\cfind{x}(y)}) =
\begin{cases}
1 & \text{si $\cfind{x}(y) \fdefined$} \\
0 & \text{sinon.}
\end{cases}
\]

Non. Sinon, la fonction partielle
\[
g(x) \simeq
\begin{cases}
0 & \text{si $h(\gn{\cfind{x}(x)}) = 0$} \\
\text{indéfinie} & \text{sinon}
\end{cases}
\]
serait calculable et aurait donc un indice~$e$.
Mais alors
\[
\cfind{e}(e) \simeq
\begin{cases}
0 & \text{si $h(\gn{\cfind{e}(e)}) = 0$} \\
\text{indéfinie} & \text{sinon,}
\end{cases}
\]
si bien que $\cfind{e}(e)$ serait définie si
et seulement si elle ne l'est pas : contradiction.

Regardons maintenant l'équation où figure
$\cfind{e}$. Elle contient une forme
d'autoréférence : la valeur de $\cfind{e}(e)$
dépend du « nom » $\gn{\cfind{e}(e)}$.
Le théorème du point fixe affirme qu'une telle
construction est possible en général, et non
seulement pour produire des contradictions.

Le \olref{lem:fixed-equiv} en donne deux
formulations équivalentes. Puisqu'elles sont
toutes deux vraies, elles sont, au sens logique,
équivalentes; ce qui nous intéresse surtout est
que chacune se déduit directement de l'autre.
Elles expriment ainsi le même théorème.

\begin{lem}
\ollabel{lem:fixed-equiv}
Les deux énoncés suivants sont équivalents :
\begin{enumerate}
\item Pour toute fonction calculable partielle
  $g(x,y)$, il existe un indice~$e$ tel que,
  pour tout~$y$,
\[
\cfind{e}(y) \simeq g(e,y).
\]
\item Pour toute fonction calculable~$f(x)$,
  il existe un indice~$e$ tel que, pour tout~$y$,
\[
\cfind{e}(y) \simeq \cfind{f(e)}(y).
\]
\end{enumerate}
\end{lem}

\begin{proof}
$(1)\Rightarrow(2)$ : étant donnée~$f$,
définissons $g$ par
$g(x,y)\simeq\fn{Un}(f(x),y)$.
L'énoncé (1) fournit un indice~$e$ tel que,
pour tout~$y$,
\begin{align*}
\cfind{e}(y) & \simeq \fn{Un}(f(e),y) \\
& \simeq \cfind{f(e)}(y).
\end{align*}

$(2)\Rightarrow(1)$ : étant donnée~$g$,
le théorème $s$-$m$-$n$ fournit une fonction récursive primitive, donc totale calculable,~$f$
telle que, pour tous $x$ et~$y$,
$\cfind{f(x)}(y)\simeq g(x,y)$.
Appliquons (2) pour obtenir un indice~$e$ tel que
\begin{align*}
\cfind{e}(y) & \simeq \cfind{f(e)}(y) \\
& \simeq g(e,y).
\end{align*}
La démonstration est achevée.
\end{proof}

\begin{explain}
Avant de prouver que (1) est vrai, et donc que
(2) l'est aussi, mesurons ce qu'il affirme.
Pensons à l'indice~$e$ comme à un programme :
pour toute fonction calculable partielle~$g(x,y)$,
on peut trouver un programme~$e$ qui calcule
$g_e(y)\simeq g(e,y)$. Ce programme calcule
donc une fonction qui fait référence au
programme lui-même.
\end{explain}

\begin{thm}
Les deux énoncés du \olref{lem:fixed-equiv}
sont vrais. En particulier, pour toute fonction
calculable partielle~$g(x,y)$, il existe un
indice~$e$ tel que, pour tout~$y$,
\[
\cfind{e}(y) \simeq g(e,y).
\]
\end{thm}

\begin{proof}
Les ingrédients de la preuve étaient déjà
présents dans la discussion du problème de
l'arrêt. Soit $\fn{diag}(x)$ une fonction
calculable qui, pour chaque~$x$, renvoie un
indice de la fonction
$f_x(y)\simeq\cfind{x}[2](x,y)$, soit
\[
\cfind{\fn{diag}(x)}(y) \simeq \cfind{x}[2](x,y).
\]
La fonction~$\fn{diag}$ transforme donc le
programme d'une fonction binaire en programme
d'une fonction unaire, en fixant le premier
argument égal au code du programme initial.
On peut la définir formellement. Posons d'abord
\[
s(x,y) \simeq \fn{Un}^2(x,x,y),
\]
où $\fn{Un}^2$ est une fonction ternaire
universelle pour les fonctions calculables
partielles binaires. Le théorème $s$-$m$-$n$
donne alors une fonction récursive primitive
$\fn{diag}$ telle que
\[
\cfind{\fn{diag}(x)}(y) \simeq s(x,y).
\]

Définissons maintenant $l$ par
\[
l(x,y) \simeq g(\fn{diag}(x),y).
\]
Soit $\gn{l}$ un indice de~$l$, et posons
$e=\fn{diag}(\gn{l})$. Pour tout~$y$, on a
\begin{align*}
\cfind{e}(y) & \simeq \cfind{\fn{diag}(\gn{l})}(y) \\
& \simeq \cfind{\gn{l}}[2](\gn{l}, y) \\
& \simeq l(\gn{l}, y) \\
& \simeq g(\fn{diag}(\gn{l}),y) \\
& \simeq g(e, y),
\end{align*}
comme voulu.\footnote{Corrections éditoriales : les évaluations diagonales portent sur des fonctions binaires, d'où les marqueurs d'arité $[2]$. Les égalités entre valeurs éventuellement indéfinies s'écrivent $\simeq$ ; la chaîne source omet ces précisions.}
\end{proof}

\begin{explain}
Que se passe-t-il concrètement ? Supposons
que l'on veuille écrire un programme qui
affiche son propre texte. Imaginons un langage
doté d'une bibliothèque de fonctions sur les
chaînes de caractères particulièrement riche.
La fonction $\fn{diag}$ y fonctionne ainsi :
à partir d'une chaîne~$s$, elle repère chaque
occurrence du symbole `x' dans~$s$ et la
remplace par une copie entre guillemets de
la chaîne d'origine. Par exemple, appliquée à
\begin{quote}
\begin{verbatim}
hello x world
\end{verbatim}
\end{quote}
elle renvoie
\begin{quote}
\begin{verbatim}
hello 'hello x world' world
\end{verbatim}
\end{quote}
Dans ce langage, le programme recherché
s'écrit simplement
\begin{quote}
\begin{verbatim}
print(diag('print(diag(x))'))
\end{verbatim}
\end{quote}
On peut vérifier qu'il produit son propre
texte. Dans des langages usuels comme C++ ou
Java, le même principe fonctionne, avec une
réalisation plus compliquée.

Il ne reste que quelques étapes avant de
retrouver la preuve du théorème du point fixe.
Supposons qu'une variante de la fonction
$\fn{print}(x,y)$ prenne une chaîne~$x$ et
un entier~$y$, puis affiche $x$ exactement
$y$ fois. Alors le « programme »
\begin{quote}
\begin{verbatim}
getinput(y);
print(diag('getinput(y); print(diag(x), y)'), y)
\end{verbatim}
\end{quote}
affiche son propre texte, à l'espacement de présentation près, $y$ fois sur
l'entrée~$y$. En remplaçant le squelette
$\fn{getinput}$--$\fn{print}$--$\fn{diag}$
par une fonction arbitraire~$g(x,y)$, on obtient
\begin{quote}
\begin{verbatim}
g(diag('g(diag(x), y)'), y)
\end{verbatim}
\end{quote}
où $y$ désigne l'entrée du programme ; cette expression applique~$g$ au
programme lui-même et à~$y$. Si l'on comprend
« citer » comme « utiliser un indice de »,
on retrouve la preuve formelle ci-dessus.

On peut pour l'instant considérer cette preuve
comme une astuce formelle, voire comme de la
magie noire. Il faut néanmoins pouvoir en
reconstruire les étapes. Lors de l'étude des
théorèmes d'incomplétude et du théorème du
point fixe qui leur est lié, nous verrons
d'autres façons de comprendre le mécanisme.
\end{explain}

\begin{tagblock}{lambda}
\begin{digress}
La même idée permet d'obtenir un combinateur
de point fixe. Soit un terme~$g$ du lambda-calcul;
on cherche un terme~$k$ qui soit $\beta$-équivalent
à~$gk$. Définissons, selon nos conventions de
notation, les termes
\[
\fn{diag}(x) = xx
\]
et
\[
l(x) = g(\fn{diag}(x))
\]
On choisit $x$ sans occurrence libre dans $g$ ; autrement dit, $l=\lambd[x][g(xx)]$. Posons
$k=ll$. On a alors
\begin{align*}
k & = (\lambd[x][g(xx)])(\lambd[x][g(xx)]) \\
& \red  g((\lambd[x][g(xx)])(\lambd[x][g(xx)])) \\
& = gk.
\end{align*}
Si l'on pose
\[
Y = \lambd[g][((\lambd[x][g(xx)])(\lambd[x][g(xx)]))]
\]
alors $Yg$ et $g(Yg)$ se réduisent à un terme
commun; donc $Yg\equiv_\beta g(Yg)$. C'est le
« combinateur de Curry ». Si, en revanche,
on pose
\[
Y = (\lambd[xg][g(xxg)])(\lambd[xg][g(xxg)])
\]
alors $Yg$ se réduit effectivement à $g(Yg)$,
ce qui est un énoncé plus fort. Cette dernière
version de~$Y$ est le « combinateur de Turing ».
\end{digress}
\end{tagblock}

\end{document}
